\section{理论原理}

\subsection{pn结势垒电容与杂质浓度的关系}

\subsubsection{单边突变结模型}

当pn结中一边的掺杂浓度远高于另一边时（如 $\mathrm{p}^+\mathrm{n}$ 结中 $N_{\mathrm{A}} \gg N_{\mathrm{D}}$），pn结交界面处杂质分布近似为突变，称为单边突变结。在单边突变结中，空间电荷区几乎全部扩展到轻掺杂一侧：
\begin{equation}
w = x_{\mathrm{n}} + x_{\mathrm{p}} \approx x_{\mathrm{n}}
\end{equation}
其中 $x_{\mathrm{p}}$、$x_{\mathrm{n}}$ 分别为p区和n区空间电荷区宽度，$w$ 为总空间电荷区宽度。此时只需单独求解轻掺杂一侧的泊松方程即可获得结的特性。

\subsubsection{空间电荷区宽度与偏压的关系}

在轻掺杂的n区（$0 \leqslant x \leqslant w$），泊松方程为：
\begin{equation}
\frac{\mathrm{d}^2 V}{\mathrm{d} x^2} = -\frac{\rho(x)}{\epsilon\epsilon_0} = -\frac{qN_{\mathrm{D}}}{\epsilon\epsilon_0}
\end{equation}
其中 $\epsilon$ 为半导体相对介电常量（硅的 $\epsilon = 11.8$），$\epsilon_0 = (36\pi \times 10^{11})^{-1}\,\mathrm{F/cm}$ 为真空电容率，$q$ 为电子电荷，$N_{\mathrm{D}}$ 为施主杂质浓度（假设均匀分布）。

积分一次得电场分布（取 $x = w$ 处电场为零）：
\begin{equation}
E(x) = -\frac{\mathrm{d}V}{\mathrm{d}x} = \frac{qN_{\mathrm{D}}}{\epsilon\epsilon_0}(w - x)
\end{equation}

再积分一次得电势分布。空间电荷区两端总电势差为自建势 $V_{\mathrm{D}}$（接触电势差）与外加反向偏压 $V_{\mathrm{R}}$ 之和：
\begin{equation}
V_{\mathrm{D}} + V_{\mathrm{R}} = \int_{0}^{w} E(x)\,\mathrm{d}x = \frac{qN_{\mathrm{D}}}{\epsilon\epsilon_0}\int_{0}^{w}(w - x)\,\mathrm{d}x = \frac{qN_{\mathrm{D}}w^2}{2\epsilon\epsilon_0}
\end{equation}

由此解得空间电荷区宽度与偏压的关系：
\begin{equation}\label{eq:w-V}
w = \left[\frac{2\epsilon\epsilon_0}{qN_{\mathrm{D}}}(V_{\mathrm{D}} + V_{\mathrm{R}})\right]^{1/2}
\end{equation}

\subsubsection{pn结势垒电容}

空间电荷区单位面积存储的电荷量为：
\begin{equation}
Q = qN_{\mathrm{D}}w
\end{equation}

当反向偏压增加 $\mathrm{d}V_{\mathrm{R}}$ 时，空间电荷区宽度增加 $\mathrm{d}w$，原来在 $\mathrm{d}w$ 层内的载流子被驱走，空间电荷增加 $\mathrm{d}Q$。pn结势垒电容定义为这一微分充放电过程的电容：
\begin{equation}
\frac{C}{A} = \frac{\mathrm{d}Q}{\mathrm{d}V_{\mathrm{R}}}
\end{equation}
其中 $A$ 为pn结的结面积。

利用式(\ref{eq:w-V})求导：
\begin{equation}
\frac{\mathrm{d}w}{\mathrm{d}V_{\mathrm{R}}} = \frac{\epsilon\epsilon_0}{qN_{\mathrm{D}}w}
\end{equation}

代入得：
\begin{equation}\label{eq:C-Area}
\frac{C}{A} = qN_{\mathrm{D}} \cdot \frac{\mathrm{d}w}{\mathrm{d}V_{\mathrm{R}}} = \frac{\epsilon\epsilon_0}{w}
\end{equation}

或将 $w$ 表达式直接代入：
\begin{equation}\label{eq:C-V}
C = A\left[\frac{q\epsilon\epsilon_0 N_{\mathrm{D}}}{2(V_{\mathrm{D}} + V_{\mathrm{R}})}\right]^{1/2}
\end{equation}

式(\ref{eq:C-Area})表明，pn结势垒电容在形式上与间距为 $w$、介电常量为 $\epsilon\epsilon_0$ 的平行板电容器相同。但二者的本质区别在于：空间电荷区宽度 $w$ 随外加偏压变化，因此势垒电容并非固定值，而是一个微分电容。

\subsubsection{$1/C^2$--$V$ 线性关系}

将式(\ref{eq:C-V})改写为：
\begin{equation}\label{eq:1C2-V}
\frac{1}{C^{2}} = \frac{2}{A^{2}q\epsilon\epsilon_0 N_{\mathrm{D}}}(V_{\mathrm{R}} + V_{\mathrm{D}})
\end{equation}

式(\ref{eq:1C2-V})表明，对于均匀掺杂的单边突变结，$1/C^{2}$ 与 $V_{\mathrm{R}}$ 成线性关系。由实验测得 $1/C^{2}$--$V_{\mathrm{R}}$ 直线的斜率 $k = 2/(A^{2}q\epsilon\epsilon_0 N_{\mathrm{D}})$，可求得掺杂浓度：
\begin{equation}
N_{\mathrm{D}} = \frac{2}{A^{2}q\epsilon\epsilon_0 \cdot k}
\end{equation}

将直线外推至 $1/C^{2} = 0$ 处，在电压轴上的截距即为接触电势差 $V_{\mathrm{D}}$（取绝对值）。

\subsubsection{任意杂质分布的 $N(w)$ 公式}

对于轻掺杂一侧杂质浓度不均匀的普遍情形，设空间电荷区边界 $w$ 处的杂质浓度为 $N(w)$。当偏压变化 $\mathrm{d}V_{\mathrm{R}}$ 引起空间电荷区宽度变化 $\mathrm{d}w$ 时，单位面积空间电荷的变化量为：
\begin{equation}\label{eq:dQ-arbitrary}
\mathrm{d}Q = qN(w)\,\mathrm{d}w
\end{equation}

由泊松方程，电荷增量引起电场变化：
\begin{equation}
\mathrm{d}E = \frac{\mathrm{d}Q}{\epsilon\epsilon_0}
\end{equation}

相应的电势变化：
\begin{equation}
\mathrm{d}(V_{\mathrm{R}} + V_{\mathrm{D}}) = \mathrm{d}V_{\mathrm{R}} = w\,\mathrm{d}E
\end{equation}

消去 $\mathrm{d}E$ 得：
\begin{equation}\label{eq:dQdV-w}
\frac{\mathrm{d}Q}{\mathrm{d}V_{\mathrm{R}}} = \frac{\epsilon\epsilon_0}{w} = \frac{C}{A}
\end{equation}

将式(\ref{eq:dQ-arbitrary})与式(\ref{eq:dQdV-w})联立：
\begin{equation}
qN(w)\,\mathrm{d}w = \frac{C}{A}\,\mathrm{d}V_{\mathrm{R}}
\end{equation}

由 $C/A = \epsilon\epsilon_0/w$ 得 $w = \epsilon\epsilon_0 A/C$，从而：
\begin{equation}
\mathrm{d}w = -\frac{\epsilon\epsilon_0 A}{C^{2}}\,\mathrm{d}C
\end{equation}

代入整理：
\begin{equation}\label{eq:Nw-C}
N(w) = -\frac{C^{3}}{q\epsilon\epsilon_0 A^{2}\,(\mathrm{d}C/\mathrm{d}V_{\mathrm{R}})}
\end{equation}

由于反向偏压下 $\mathrm{d}C/\mathrm{d}V_{\mathrm{R}} < 0$，实际计算中使用绝对值形式：
\begin{equation}\label{eq:Nw-C-abs}
N(w) = \frac{C^{3}}{q\epsilon\epsilon_0 A^{2}\left|\mathrm{d}C/\mathrm{d}V_{\mathrm{R}}\right|}
\end{equation}

亦可用 $1/C^{2}$ 的导数表达。由 $\mathrm{d}(1/C^{2})/\mathrm{d}V_{\mathrm{R}} = -2(\mathrm{d}C/\mathrm{d}V_{\mathrm{R}})/C^{3}$，得：
\begin{equation}\label{eq:Nw-1C2}
N(w) = \frac{2}{q\epsilon\epsilon_0 A^{2}\,\dfrac{\mathrm{d}(1/C^{2})}{\mathrm{d}V_{\mathrm{R}}}}
\end{equation}

式(\ref{eq:Nw-1C2})更为实用：从实验 $C$--$V$ 数据作出 $1/C^{2}$--$V_{\mathrm{R}}$ 曲线，求各偏压下的斜率即可得到 $N(w)$。同时，各偏压下对应的空间电荷区宽度由 $w = \epsilon\epsilon_0 A/C$ 确定，从而得到 $N(w)$--$w$ 杂质分布曲线。


\subsection{势垒电容的测量原理}

pn结势垒电容是随直流偏压变化的微分电容，测量时需在pn结上同时施加直流反向偏压 $V_{\mathrm{R}}$ 和微小的交流小信号 $v(t)$。

\subsubsection{测量线路}

测量线路如图\ref{fig:measure-circuit}所示（示意图，详见实验讲义图9-2-10）。待测pn结电容 $C_{\mathrm{x}}$ 与已知电容 $C_{0}$ 串联，交流信号 $v(t)$ 加在串联电路两端，$C_{0}$ 两端电压 $v_{\mathrm{i}}$ 作为输出电压送入锁定放大器。

\subsubsection{$C_{0} \gg C_{\mathrm{x}}$ 条件下的近似}

两个电容串联，$C_{0}$ 两端的分压为：
\begin{equation}
v_{\mathrm{i}} = v(t) \cdot \frac{\dfrac{1}{\mathrm{j}\omega C_{0}}}{\dfrac{1}{\mathrm{j}\omega C_{\mathrm{x}}} + \dfrac{1}{\mathrm{j}\omega C_{0}}}
= v(t) \cdot \frac{C_{\mathrm{x}}}{C_{\mathrm{x}} + C_{0}}
\end{equation}

当 $C_{0} \gg C_{\mathrm{x}}$ 时，$C_{\mathrm{x}} + C_{0} \approx C_{0}$，得：
\begin{equation}\label{eq:meas-principle}
v_{\mathrm{i}} \approx \frac{C_{\mathrm{x}}}{C_{0}}\,v(t)
\end{equation}

式(\ref{eq:meas-principle})表明，$C_{0}$ 上的交变电压 $v_{\mathrm{i}}$ 与待测电容 $C_{\mathrm{x}}$ 成正比，比例系数为 $v(t)/C_{0}$。通过锁定放大器精确测定 $v_{\mathrm{i}}$ 的幅值，即可获得 $C_{\mathrm{x}}$。

\subsubsection{测量条件约束}

\begin{enumerate}
  \item $C_{0} \gg C_{\mathrm{x}}$：保证近似条件成立。本实验中 $C_{0}$ 选取 $5 \times 10^{3}\,\mathrm{pF}$，待测pn结零偏电容约几十 $\mathrm{pF}$。
  \item $v(t) \ll V_{\mathrm{D}} + V_{\mathrm{R}}$：交流小信号远小于pn结上的总直流压降，保证微分电容定义的线性条件。本实验 $v(t)$ 选取约 $40\,\mathrm{mV}$，硅pn结内建电势约 $0.6$--$0.8\,\mathrm{V}$，满足条件。
  \item 交流信号频率应避开 $1/f$ 噪声区且保证 $R_{2} \gg 1/(\omega C_{\mathrm{x}})$，防止交流信号被直流偏置电阻 $R_{2}$ 旁路。
\end{enumerate}


\subsection{锁定放大器的基本原理}

锁定放大器（Lock-in Amplifier, LIA）是一种能精确测量深埋于噪声之中周期重复信号幅值和相位的交流电压表。其核心由相敏检波器（PSD）和低通滤波器（LPF）组成，二者合称为相关器。

\subsubsection{相敏检波器（PSD）}

相敏检波器本质上是一个模拟乘法器，其输出电压 $v_{0}$ 是输入信号 $v_{\mathrm{i}}$ 与参考信号 $v_{\mathrm{R}}$ 的乘积。设二者均为正弦波：
\begin{align}
v_{\mathrm{i}} &= E_{\mathrm{i}}\sin(\omega_{1}t + \phi_{1}) \\
v_{\mathrm{R}} &= E_{\mathrm{R}}\sin(\omega_{2}t + \phi_{2})
\end{align}

则输出电压为：
\begin{equation}\label{eq:PSD-output}
v_{0} = v_{\mathrm{i}} \cdot v_{\mathrm{R}} = \frac{E_{\mathrm{i}}E_{\mathrm{R}}}{2}\cos\big[(\omega_{1}-\omega_{2})t + (\phi_{1}-\phi_{2})\big]
- \frac{E_{\mathrm{i}}E_{\mathrm{R}}}{2}\cos\big[(\omega_{1}+\omega_{2})t + (\phi_{1}+\phi_{2})\big]
\end{equation}

式(\ref{eq:PSD-output})包含差频分量 $(\omega_{1} - \omega_{2})$ 与和频分量 $(\omega_{1} + \omega_{2})$。当 $\omega_{1} = \omega_{2}$（信号与参考同频）时，差频分量为直流：
\begin{equation}
v_{0,\mathrm{DC}} = \frac{E_{\mathrm{i}}E_{\mathrm{R}}}{2}\cos(\phi_{1} - \phi_{2})
\end{equation}

其幅值正比于输入信号幅值 $E_{\mathrm{i}}$，且与相位差 $(\phi_{1} - \phi_{2})$ 的余弦成正比——这正是"相敏"的含义。通过调节参考通道中的移相器使 $\phi_{1} - \phi_{2} = 0$，可获得最大直流输出。

实际锁定放大器中常采用开关型PSD，参考信号为方波。方波幅度稳定，其傅里叶展开为：
\begin{equation}
v_{\mathrm{R}} = \frac{4}{\pi}\sum_{n=0}^{\infty}\frac{1}{2n+1}\sin\big[(2n+1)(2\pi f_{2}t + \phi_{2})\big]
\end{equation}

开关型PSD对参考方波的所有奇次谐波均产生相敏响应，响应幅度与谐波次数 $(2n+1)$ 成反比。

\subsubsection{低通滤波器（LPF）}

和频分量以及频率不同于参考频率的噪声分量需由低通滤波器滤除。典型的有源一阶低通滤波器传递函数为：
\begin{equation}
\frac{\dot{u}_{\mathrm{sc}}}{\dot{u}_{\mathrm{sr}}} = -K\frac{1}{1 + \mathrm{j}\omega RC}
\end{equation}

其等效噪声带宽为：
\begin{equation}\label{eq:noise-bandwidth}
\Delta f_{\mathrm{N}} = \int_{0}^{\infty}\left|\frac{1}{1 + \mathrm{j}\omega RC}\right|^{2}\mathrm{d}f = \frac{1}{4RC}
\end{equation}

时间常量 $T = RC$ 越大，等效噪声带宽越窄。对于白噪声，噪声电压正比于 $\sqrt{\Delta f_{\mathrm{N}}}$，因此锁定放大器对信噪比的改善为：
\begin{equation}
\frac{\text{输出信噪比}}{\text{输入信噪比}} = \sqrt{\frac{\Delta f_{\text{输入}}}{\Delta f_{\text{输出}}}}
\end{equation}

例如，输入噪声带宽 $10\,\mathrm{kHz}$、$T = 1\,\mathrm{s}$ 时，$\Delta f_{\text{输出}} = 0.25\,\mathrm{Hz}$，信噪比改善 $200$ 倍（$46\,\mathrm{dB}$）。

\subsubsection{相关器}

相敏检波器与低通滤波器的组合构成相关器。其输出为输入信号与参考信号的相关函数：
\begin{equation}
V_{0}' = \lim_{T \to \infty}\frac{1}{2T}\int_{-T}^{T} V_{\mathrm{i}}(t) \cdot V_{\mathrm{R}}(t - \tau)\,\mathrm{d}t
\end{equation}
其中 $\tau = (\phi_{1} - \phi_{2})/\omega$ 为参考信号相对于输入信号的延迟时间。

考虑输入信号为正弦波与噪声的叠加：
\begin{equation}
V_{\mathrm{i}} = E_{\mathrm{s}}\cos(\omega_{0}t + \theta) + E_{\mathrm{n}}\cos(\omega t + \alpha),\quad V_{\mathrm{R}} = E_{\mathrm{R}}\cos\omega_{0}t
\end{equation}

经过PSD后的输出包含直流项、$2\omega_{0}$ 和频项、$(\omega \pm \omega_{0})$ 差频与和频项。经LPF后，高频分量均被滤除，仅保留：
\begin{equation}
V_{0}' = \frac{1}{2}E_{\mathrm{s}}E_{\mathrm{R}}\cos\theta + \frac{1}{2}E_{\mathrm{n}}E_{\mathrm{R}}\cos(\Delta\omega t + \alpha)
\end{equation}
其中 $\Delta\omega = \omega - \omega_{0}$。上式第一项为与信号幅值成正比的直流输出；第二项为噪声贡献，仅当 $|\Delta\omega|$ 落在LPF等效噪声带宽内时才对输出有影响。

相关器由此起到"梳齿滤波器"的作用：在参考频率 $f_{2}$ 及其所有奇次谐波处，各有一个带宽为 $\Delta f_{\mathrm{N}} = 1/(4RC)$ 的窄带通滤波器。与参考信号相干的信号全部通过，而宽带噪声的绝大部分被抑制。一个深埋于噪声之中（如输入信噪比 $1:10$）的微弱信号，经锁定放大器后可被可靠检测（输出信噪比 $20:1$）。


\subsection{利用C--V曲线求杂质浓度分布的方法}

综合以上理论，从实验 $C$--$V$ 数据求取杂质浓度分布 $N(w)$ 的步骤为：

\begin{enumerate}
  \item \textbf{获取 $C$--$V_{\mathrm{R}}$ 数据}：在不同反向偏压 $V_{\mathrm{R}}$ 下，由锁定放大器输出 $V_{0}' \propto C_{\mathrm{x}}$，结合标准电容标定得到各偏压下的势垒电容 $C$。

  \item \textbf{计算各偏压下的空间电荷区宽度}：
  \begin{equation}
  w(V_{\mathrm{R}}) = \frac{\epsilon\epsilon_0 A}{C(V_{\mathrm{R}})}
  \end{equation}

  \item \textbf{计算杂质浓度 $N(w)$}：

  \textit{方法一}——利用 $1/C^{2}$--$V_{\mathrm{R}}$ 曲线的斜率：
  \begin{equation}
  N(w) = \frac{2}{q\epsilon\epsilon_0 A^{2}\,\dfrac{\mathrm{d}(1/C^{2})}{\mathrm{d}V_{\mathrm{R}}}}
  \end{equation}

  \textit{方法二}——利用 $C$--$V_{\mathrm{R}}$ 曲线的斜率：
  \begin{equation}
  N(w) = \frac{C^{3}}{q\epsilon\epsilon_0 A^{2}\left|\mathrm{d}C/\mathrm{d}V_{\mathrm{R}}\right|}
  \end{equation}

  \item \textbf{绘制 $N(w)$--$w$ 曲线}：以 $w$ 为横坐标、$N(w)$ 为纵坐标，即得轻掺杂一侧的杂质浓度分布。

  \item \textbf{外推求接触电势差 $V_{\mathrm{D}}$}：对于均匀掺杂段，$1/C^{2}$--$V_{\mathrm{R}}$ 为直线，外推至 $1/C^{2}=0$ 处，电压轴截距即为 $-V_{\mathrm{D}}$。
\end{enumerate}
